音声認識エンジンの実用化に取り組み、深層学習言語モデルを用いた認識性能向上と製品化を進めました。Linux、shell、Python、Git、Chainer、C/C++ などを用いて研究開発を行い、複数の関連特許や RECAIUS 音声認識サービスにつながりました。

また、会議ドメインに特化した音声認識モデルの構築にも取り組みました。会議データの収集や整備を進め、実運用に近い条件で性能を高めました。

さらに、少量データからの言語モデル適応技術にも取り組みました。限られたデータからドメイン適応を行うために、テキスト抽出、情報検索、クラスタリング、トピックモデルなどの技術を活用し、実用的な音声認識モデルの構築につなげました。
